\chapter{Metodologia Proposta}
\label{cap:metodologia}

% Descrever o pipeline completo, mesmo que parcialmente implementado.

\section{Visão geral do pipeline}
\label{sec:visao-geral-pipeline}

% Tópicos sugeridos:
%   - diagrama/figura das etapas do pipeline
%   - o que já está implementado x o que é proposto/planejado

\section{Geração de fragmentos e combinação}
\label{sec:etapa-fragmentos}

% Tópicos sugeridos:
%   - categorias de fragmentos utilizadas (anéis, auxocromos, pontes, linhas)
%   - regras de combinação/recombinação
%   - validação química e deduplicação

\section{Representação estrutural (árvore/grafo)}
\label{sec:representacao-estrutural}

% Tópicos sugeridos:
%   - como a molécula (SMILES) é convertida em árvore/grafo
%   - papel do algoritmo de Zhang-Shasha / GA-grafo nessa representação

\section{Avaliação de diversidade estrutural}
\label{sec:avaliacao-diversidade}

% Tópicos sugeridos:
%   - métrica proposta (distância de edição de árvores)
%   - como será usada: caracterização do conjunto gerado, filtragem de
%     redundância, comparação com baseline

\section{Ambiente e ferramentas}
\label{sec:ambiente-ferramentas}

% Tópicos sugeridos:
%   - linguagem, bibliotecas (RDKit, pandas etc.), versões
